\chapter[INTRODUCTION]{Introduction}

This chapter introduces the financial decision-support problem behind TrustFin AI. It explains why a loan prediction needs more context than an approval label, sets out the project's aim and objectives, and defines the scope of the application. The final part outlines how the remaining chapters develop the design, implementation and evaluation.

\section{Background and Motivation}
\subsection{Project Context and Credit Assessment}
A loan application brings together several facts that have to be considered together: the amount requested, the applicant's income, existing monthly repayments, credit information and the length of the proposed loan. A high income may appear reassuring until the existing debt is considered. Similarly, a long repayment term can reduce the monthly instalment while increasing the period over which a borrower remains committed. A useful assessment needs to preserve these relationships instead of treating each field as an isolated number.

TrustFin AI is a final-year project that explores this problem through two loan products: a residential home loan and an unsecured personal loan. Each product has a separate XGBoost classifier and preprocessing pipeline. The application accepts an applicant profile, returns an approval or rejection prediction, and presents the probability assigned by the model. It then produces two local explanations, using SHAP and LIME, and calculates four diagnostic scores.

The project goes beyond a notebook that prints a classification result. It includes a browser interface, a FastAPI service, stored evaluation history, batch processing and PDF reporting. These parts make it possible to follow an application from submitted inputs to its saved assessment. The datasets, however, are generated for the project. They provide a controlled setting for development rather than evidence about actual borrowers.

Credit scoring has long used statistical classification to support lending decisions. Hand and Henley's review discusses consumer credit as a specialised modelling problem, while later benchmarking work shows the value of comparing classifiers on relevant credit datasets \cite{handHenley1997,lessmann2015}. The lesson for this project is practical: an algorithm's reputation does not replace evaluation on the data and target that it will actually use.

A small scorecard can often be explained by showing a few rules and their weights. A boosted ensemble combines many trees, and its overall behaviour is harder to communicate through one short rule. XGBoost is useful for mixed numerical and categorical records after preprocessing, but a user still needs a way to inspect an individual output. The prediction and its explanation therefore serve different purposes.

In TrustFin AI, the target label is simulated underwriting approval. It is not an observed default or repayment event. This distinction matters whenever the interface displays a probability. A model trained on approval labels estimates the approval class under those labels; it has not learned a measured probability that a borrower will repay a loan.

\subsection{Need for Explainable Predictions}
An applicant who sees only ``Rejected'' has little information to work with. A reviewing officer may have the same difficulty when the model's result conflicts with a simple reading of the application. Local explanations can identify the features associated with that particular output and give both users a starting point for review.

SHAP and LIME provide two different views of the same model response. SHAP assigns feature contributions relative to an explanation baseline. LIME builds a simpler model around the selected input. Their outputs can be displayed as ranked feature contributions, but the methods do not use identical assumptions or output units \cite{lundberg2017,ribeiro2016}. Agreement is useful to inspect; disagreement needs interpretation rather than automatic rejection of one method.

An explanation is also not a guarantee that the model behaves well. A model may rely on an unsuitable attribute, and an accurate explanation may reveal that dependence. Conversely, a plausible-looking explanation may describe the model poorly. The project therefore includes checks on the explanations instead of treating a generated chart as sufficient evidence of reliability.

\subsection{Project Motivation}
The initial question that led to the creation of TrustFin AI is what extra information should be included in the output of a loan model, so that a person can sensibly review it? The application does so by addressing the following local reasons, repayment calculations, explanation diagnostics and saved assessment records.

The two-product structure provides detail to the problem that is useful. While property value and loan-to-value ratio have been important to the home-loan pipeline, purpose and credit-card utilisation are important to the personal-loan pipeline. Having separate models prevents interpretation of fields that are not relevant to the product.

The second reason is that it helps to maintain separate classifications of various types of evidence. Predictive accuracy is about the performance on labelled records. The output probability is known as model confidence. An explanation is a description of the contributions in features and the Trust Score is a combination of the project's diagnostic checks. These values are all present in the workflow, but should not be interpreted as being synonymous with measures of quality.

\section{Aim and Objectives}
The goal is to create an application for decision-making support, where the prediction is product-specific and there are local explanations for this prediction, as well as explicit checks of these explanations.

The aims are to:
\begin{enumerate}
\item Create two distinct datasets of home loans and personal loans;
\item  Train XGBoost classifiers using a fixed, stratified train--test split;
\item Derive repayment features, and consistently use the fitted preprocessing during inference;
\item Compute local SHAP and LIME contributions for an individual application;
\item Conduct the faithfulness, stability and consistency and Gender-sensitivity diagnostics;
\item Combine the three explanation diagnostics in the Trust Score while reporting Gender sensitivity and prediction probability separately;
\item Present, store, retrieve and export assessments via the web application.
\end{enumerate}

The current software are described by these objectives. They are not relying on the validity of the resulting model or Trust Score for lending operations to have been independently validated.

\section{Scope of the Project}
The application supports single-record evaluation, product-specific forms, batch CSV and JSON processing, an applicant view, bank-officer views, evaluation history and downloadable PDF reports. A what-if calculator allows users to change selected affordability inputs and inspect the resulting instalment and repayment ratio.

The calculator performs financial arithmetic in the browser. Moving a slider does not submit a fresh prediction request. Its result is therefore an affordability illustration, not a promise that the model will approve a changed application. The officer and applicant modes are also presentation choices; the source does not implement authenticated accounts or server-side role permissions.

In the current work, there is no link to a bank, no retrieval of real time bureau information, there is no setting of an institution's credit policy nor a legally binding lending decision. However, repayment rates information and probability calibration, secure deployment and population fairness monitoring are not in the scope of the evaluated. The boundaries enable the project to be evaluated on the project's actual delivery - an integrated academic system for inspection of model based loan assessments.

\section{Approach Followed}
The process starts when the user chooses a product and fills a form requesting information about the product. The backend normalises the values sent, computes financial features and applies the selected preprocessor to the value. XGBoost provides the class and approval probability. SHAP and LIME then give explanations for the modified record and the validation functions compute their diagnostic values.

The output of the response is separated, thereby allowing for a non-collapsing of the interface showing the model result, solvency information and explanation checks. The route also will try to preserve the assessment. The response structure is the same for History and PDF, making it easier for the user to follow a displayed response to the submitted application.

\section{Organisation of the Report}
Chapter 2 reviews related research and Chapter 3 sets out the requirements. Design and Implementation are discussed in Chapters 4 and 5. Verification and results are covered in Chapters 6 and 7. The work plan is presented in Chapter 8 and limitations and future directions are given in Chapter 9.

\par\medskip\noindent
This chapter has established the purpose of TrustFin AI, the roles of prediction and explanation, and the limits of the application. It has also defined the objectives that guide the technical chapters that follow.
